\documentclass[12pt,a4paper]{article}

\usepackage[margin=1in]{geometry}
\usepackage[T1]{fontenc}
\usepackage[utf8]{inputenc}
\usepackage{lmodern}
\usepackage{setspace}
\usepackage{hyperref}
\usepackage{enumitem}
\usepackage{longtable}
\usepackage{booktabs}
\usepackage{xcolor}

\hypersetup{
  colorlinks=true,
  linkcolor=black,
  urlcolor=blue,
  pdftitle={CareMate Comprehensive Technical Project Report},
  pdfauthor={CareMate Team}
}

\setstretch{1.15}
\setlist[itemize]{leftmargin=1.5em}
\setlist[enumerate]{leftmargin=1.7em}

\title{\textbf{CareMate: Comprehensive Technical Project Report}\\\large(Current Working Model)}
\author{CareMate Project Team}
\date{Codebase Snapshot: March 30, 2026}

\begin{document}

\maketitle
\tableofcontents
\newpage

\section{Project Snapshot}
\begin{itemize}
  \item \textbf{Project:} CareMate (assistive app for Alzheimer's care and memory support)
  \item \textbf{Architecture style:} Client-server web application with AI-assisted processing and local persistence
  \item \textbf{Primary runtime split:}
  \begin{itemize}
    \item Backend: Python + FastAPI + SQLite + ML/AI services
    \item Frontend: React + Vite single-page application
  \end{itemize}
\end{itemize}

This report is based on the currently implemented source code and active environment configuration in the project workspace.

\section{Current Working AI Model Configuration}
From the active backend environment (\texttt{CAREMATE/.env}), the configured AI setup is:
\begin{itemize}
  \item \texttt{CARE\_LLM\_PROVIDER=gemini}
  \item \texttt{CARE\_LLM\_MODEL=gemini-2.5-flash-lite}
  \item \texttt{CARE\_TRANSCRIBE\_PROVIDER=gemini}
  \item \texttt{CARE\_TRANSCRIBE\_MODEL=gemini-2.5-flash-lite}
\end{itemize}

\textbf{Implications:}
\begin{itemize}
  \item Conversation transcription is attempted through Gemini first.
  \item Summary and reminder extraction are attempted through Gemini first.
  \item If AI transcription fails or is unavailable, backend falls back to local Whisper (\texttt{base}).
  \item If AI summary extraction fails or is unavailable, backend falls back to local BART summarization plus heuristic reminder extraction.
\end{itemize}

\section{High-Level Technical Architecture}

\subsection{Backend Layer (\texttt{CAREMATE/})}
\textbf{Core files:}
\begin{itemize}
  \item \texttt{main.py} (832 lines): API routes, face pipeline, audio pipeline, reminders, diary CRUD
  \item \texttt{llm.py} (230 lines): provider abstraction for Gemini/OpenAI-compatible endpoints
  \item \texttt{database.py} (27 lines): SQLite schema creation
  \item \texttt{init\_db.py}: database bootstrap script
\end{itemize}

\textbf{Core backend capabilities:}
\begin{itemize}
  \item Audio ingestion and conversion to WAV via FFmpeg
  \item AI-first transcription and summarization flow with local fallback
  \item Face registration, identification, and verification using DeepFace
  \item Reminder extraction/storage and status lifecycle management
  \item Diary entry persistence and retrieval
\end{itemize}

\subsection{Frontend Layer (\texttt{caremate\_front/})}
\textbf{Core files:}
\begin{itemize}
  \item \texttt{src/App.jsx} (700 lines): end-to-end app flow in a single module
  \item \texttt{src/main.jsx}: React root initialization
  \item \texttt{src/App.css}, \texttt{src/index.css}: design system and layout styles
\end{itemize}

\textbf{Core frontend capabilities:}
\begin{itemize}
  \item Camera and microphone capture flow
  \item Face recognition-assisted person identification
  \item Conversation recording and backend processing trigger
  \item Diary browsing and deletion
  \item Face library management (add/delete/select/rename via API)
  \item Reminder list management and completion toggle
\end{itemize}

\section{Backend Technology Stack}

\subsection{Framework and Runtime}
\begin{itemize}
  \item FastAPI (\texttt{fastapi==0.129.0})
  \item Uvicorn (\texttt{uvicorn==0.41.0})
  \item CORS middleware enabled (currently wildcard \texttt{allow\_origins=["*"]})
  \item Multipart upload handling via \texttt{python-multipart}
\end{itemize}

\subsection{AI and ML}
\begin{itemize}
  \item Local STT fallback: \texttt{openai-whisper} using model \texttt{base}
  \item Local summarization fallback: \texttt{facebook/bart-large-cnn} via Transformers
  \item Face recognition: DeepFace with ArcFace model and MTCNN detector backend
  \item Cloud AI support:
  \begin{itemize}
    \item Gemini API (\texttt{generativelanguage.googleapis.com})
    \item OpenAI-compatible providers via configurable base URL
  \end{itemize}
\end{itemize}

\subsection{Data and Utilities}
\begin{itemize}
  \item SQLite for persistence (\texttt{caremate.db})
  \item \texttt{python-dateutil} for due-date parsing
  \item \texttt{requests} for provider API integration
  \item FFmpeg subprocess for audio normalization (16kHz mono WAV)
\end{itemize}

\section{Frontend Technology Stack}
\begin{itemize}
  \item React 19
  \item Vite 8
  \item Axios for HTTP calls
  \item Framer Motion for transitions and animated interaction states
  \item Lucide React icon system
  \item CSS variable-driven visual system
\end{itemize}

\textbf{Build and dev scripts:}
\begin{itemize}
  \item \texttt{npm run dev} --- Vite development server
  \item \texttt{npm run build} --- production bundle
  \item \texttt{npm run preview} --- built app preview
\end{itemize}

\section{Data Model (SQLite)}

\subsection{\texttt{conversations}}
\begin{itemize}
  \item \texttt{id} INTEGER PRIMARY KEY AUTOINCREMENT
  \item \texttt{person\_name} TEXT
  \item \texttt{transcript} TEXT
  \item \texttt{summary} TEXT
  \item \texttt{created\_at} TIMESTAMP DEFAULT CURRENT\_TIMESTAMP
\end{itemize}

\subsection{\texttt{reminders}}
\begin{itemize}
  \item \texttt{id} INTEGER PRIMARY KEY AUTOINCREMENT
  \item \texttt{title} TEXT NOT NULL
  \item \texttt{due\_at} TIMESTAMP (nullable)
  \item \texttt{status} TEXT NOT NULL DEFAULT \texttt{pending}
  \item \texttt{conversation\_id} INTEGER (foreign key to \texttt{conversations.id})
  \item \texttt{created\_at} TIMESTAMP DEFAULT CURRENT\_TIMESTAMP
\end{itemize}

\textbf{Lifecycle:} \texttt{init\_db.py} initializes schema via \texttt{database.create\_tables()}.

\section{API Surface (Implemented)}

\subsection{Audio and Diary APIs}
\begin{longtable}{p{0.28\textwidth} p{0.67\textwidth}}
\toprule
\textbf{Endpoint} & \textbf{Purpose} \\
\midrule
\endhead
\texttt{POST /process-audio} & Accepts audio + \texttt{person\_name}; converts audio; performs AI-first transcription and analysis with fallback; stores conversation; returns saved entry and \texttt{suggested\_reminders}. \\
\texttt{GET /diary-entries} & Returns conversations ordered by newest first. \\
\texttt{DELETE /diary-entries/\{entry\_id\}} & Deletes a diary entry by ID. \\
\bottomrule
\end{longtable}

\subsection{Reminder APIs}
\begin{longtable}{p{0.28\textwidth} p{0.67\textwidth}}
\toprule
\textbf{Endpoint} & \textbf{Purpose} \\
\midrule
\endhead
\texttt{GET /reminders} & Returns reminders ordered by newest first. \\
\texttt{POST /reminders} & Creates reminder (title, optional due date, optional conversation reference). \\
\texttt{PATCH /reminders/\{reminder\_id\}} & Updates reminder status in \texttt{\{pending, done\}}. \\
\bottomrule
\end{longtable}

\subsection{Face Library and Recognition APIs}
\begin{longtable}{p{0.36\textwidth} p{0.59\textwidth}}
\toprule
\textbf{Endpoint} & \textbf{Purpose} \\
\midrule
\endhead
\texttt{GET /faces} & Lists known face identities from \texttt{known/}. \\
\texttt{GET /faces/\{person\_name\}/image} & Serves the stored face image. \\
\texttt{PATCH /faces/\{person\_name\}} & Renames a face identity. \\
\texttt{DELETE /faces/\{person\_name\}} & Deletes a face identity image. \\
\texttt{PUT /faces/\{person\_name\}/photo} & Replaces face image for an identity. \\
\texttt{POST /register-face/\{person\_name\}} & Registers an uploaded face image. \\
\texttt{POST /verify-face/\{person\_name\}} & One-to-one verification against a known identity. \\
\texttt{POST /identify-face} & One-to-many face identification against all known faces. \\
\texttt{POST /webcam-verify} & Verifies base64 webcam frame against provided person name. \\
\texttt{GET /webcam} & Serves a lightweight HTML webcam diagnostic page. \\
\bottomrule
\end{longtable}

\subsection{Contacts API}
\begin{itemize}
  \item \texttt{GET /contacts}: currently returns mocked contact data (not DB-backed).
\end{itemize}

\section{End-to-End Runtime Flows}

\subsection{Memory Capture Flow}
\begin{enumerate}
  \item Frontend opens camera/microphone via browser media APIs.
  \item Frontend captures still image and calls \texttt{/identify-face}.
  \item If no verified match exists, frontend auto-registers unknown face with generated temporary name.
  \item User records audio; frontend posts to \texttt{/process-audio}.
  \item Backend runs audio conversion, transcription, summary/reminder extraction, and DB write.
  \item Frontend displays generated summary and optional reminder suggestions.
  \item Reminder is persisted only after explicit user confirmation (Add action).
\end{enumerate}

\subsection{Reminder Management Flow}
\begin{enumerate}
  \item Frontend loads reminders from \texttt{/reminders}.
  \item Pending vs done groups are computed client-side.
  \item Toggle action sends \texttt{PATCH /reminders/\{id\}} with updated status.
\end{enumerate}

\subsection{Diary Flow}
\begin{enumerate}
  \item Frontend loads entries from \texttt{/diary-entries}.
  \item List and detail views are rendered client-side.
  \item Delete action sends \texttt{DELETE /diary-entries/\{id\}}.
\end{enumerate}

\section{Current Implementation Characteristics}

\subsection{Strengths}
\begin{itemize}
  \item AI-first with deterministic local fallback improves robustness.
  \item Clear separation between suggested reminders and persisted reminders.
  \item Face CRUD plus recognition APIs are fully connected to UI workflows.
  \item Lightweight deployment path with SQLite and a single FastAPI service.
\end{itemize}

\subsection{Technical Debt and Risks}
\begin{itemize}
  \item \texttt{main.py} and \texttt{App.jsx} are monolithic and should be modularized.
  \item CORS is fully open and not production hardened.
  \item No authentication/authorization layer for API routes.
  \item No automated backend/frontend test suite.
  \item Stateful artifact \texttt{caremate.db} exists in source tree.
  \item \texttt{/contacts} is mocked and not integrated with persistence.
  \item Broad \texttt{except Exception} handling in multiple endpoints.
  \item Heavy ML dependency footprint (TensorFlow + Torch + DeepFace).
\end{itemize}

\subsection{Performance Considerations}
\begin{itemize}
  \item Lazy model loading reduces startup cost but first request is still expensive.
  \item Face identification is linear across registered faces.
  \item Audio processing currently runs synchronously per request and writes temporary files locally.
\end{itemize}

\section{Security and Compliance Posture (Current State)}
\begin{itemize}
  \item Environment-based secrets are supported via \texttt{.env}.
  \item SQL injection risk is reduced through parameterized SQLite queries.
  \item Missing controls before production usage:
  \begin{itemize}
    \item Authentication and authorization
    \item Rate limiting and abuse controls
    \item Strict request/file size validation
    \item Audit logging and retention policy enforcement
    \item Explicit encryption-at-rest strategy beyond host defaults
  \end{itemize}
\end{itemize}

Given healthcare-adjacent usage, privacy and compliance hardening is strongly recommended before broader deployment.

\section{Deployment and Operations Notes}

\subsection{Backend Startup}
\begin{enumerate}
  \item \texttt{cd CAREMATE}
  \item Create and activate a virtual environment
  \item \texttt{pip install -r requirements.txt}
  \item \texttt{python init\_db.py}
  \item \texttt{uvicorn main:app --reload}
\end{enumerate}

\subsection{Frontend Startup}
\begin{enumerate}
  \item \texttt{cd caremate\_front}
  \item \texttt{npm install}
  \item \texttt{npm run dev}
\end{enumerate}

\subsection{Runtime Configuration Knobs}
\begin{itemize}
  \item Frontend API base: \texttt{VITE\_API\_BASE\_URL}
  \item AI provider/model config: \texttt{CARE\_LLM\_*}, \texttt{CARE\_TRANSCRIBE\_*}
  \item External requirement: \texttt{ffmpeg} available on host PATH for \texttt{/process-audio}
\end{itemize}

\section{Suggested Technical Roadmap (Priority Order)}
\begin{enumerate}
  \item \textbf{Refactor monoliths:}
  \begin{itemize}
    \item Split backend by domain (audio, faces, reminders, diary)
    \item Split frontend screens/components into independent modules
  \end{itemize}
  \item \textbf{Add test coverage:}
  \begin{itemize}
    \item Backend API tests (FastAPI TestClient)
    \item Frontend component/integration tests
  \end{itemize}
  \item \textbf{Security hardening:}
  \begin{itemize}
    \item Token-based authentication
    \item Restricted CORS origin policy
    \item Input and file validation guardrails
  \end{itemize}
  \item \textbf{Reliability improvements:}
  \begin{itemize}
    \item Structured logging with request correlation IDs
    \item Better exception taxonomy and clearer API error contracts
  \end{itemize}
  \item \textbf{Scalability improvements:}
  \begin{itemize}
    \item Background task queue for long-running audio/AI jobs
    \item Caching/indexing strategy for face matching
  \end{itemize}
  \item \textbf{Data governance:}
  \begin{itemize}
    \item Retention policy, archival process, and encrypted backup path
  \end{itemize}
\end{enumerate}

\section{Final Technical Summary}
CareMate currently operates as a practical AI-enabled caregiving assistant with a functioning FastAPI backend and React frontend. The current runtime path is Gemini-first for both transcription and conversation understanding, with local fallback logic for resiliency. The project has strong MVP-level functional depth, especially in face-aware memory capture and reminder workflows, and is ready for focused work on modularization, testing, and production-grade security controls.

\end{document}
